\documentclass[11pt]{article}
\usepackage[margin=1in]{geometry}
\usepackage{amsmath,amssymb,amsthm,booktabs,array}
\usepackage[hidelinks]{hyperref}
\newcommand{\C}{\mathbb{C}}
\newcommand{\Mil}{\operatorname{Mil}(F)}
\newtheorem{theorem}{Theorem}[section]
\newtheorem{lemma}[theorem]{Lemma}
\newtheorem{corollary}[theorem]{Corollary}
\theoremstyle{remark}
\newtheorem{remark}[theorem]{Remark}
\newtheorem{example}[theorem]{Example}

\title{On a statement of A.\,N.~Shoshitaishvili\\ about weighted homogeneous singularities}
\author{Alex Shoshitaishvili\thanks{Also published as A.\,N.~Shoshitaishvili.}}
\date{October 2026\\[4pt]{\small DOI: \href{https://doi.org/10.5281/zenodo.23405205}{10.5281/zenodo.23405205}\quad \url{https://github.com/alexshosh/statement1-weighted-homogeneous}}}

\begin{document}
\maketitle

\begin{abstract}
Let $F$ be a weighted homogeneous polynomial with an isolated critical point at $0$, with weights $w_1\ge\dots\ge w_n$ and degree $d$. Statement~1 of the author's 1976 paper says that if $\varphi(0)=0$ and the Jacobian ideal of $\varphi$ lies in the Jacobian ideal of $F$, then $\varphi$ lies in the Jacobian ideal of $F$ and has weighted degree at least $d$. The second part is equivalent to the absence of nonzero derivations of negative weight of the Milnor algebra of $F$; since 1995 Shoshitaishvili's Statement~1 has been studied as a conjecture of S.\,S.-T.~Yau, and as the Milnor-algebra case of Halperin's conjecture. We describe the history of the statement and give a proof of the following: with $M=\min(w_3-2w_n,\ w_1-3w_n)$, the Milnor algebra has no nonzero derivation of weight $-q$ if $q>M$, or if $q>M-(w_{n-1}-w_n)$ and $w_n$ does not divide $d-q$; in particular it has no nonzero derivation of negative weight when $w_3\le2w_n$ or $w_1\le3w_n$.
\end{abstract}

\section{Introduction}

\paragraph{Notation.} $P=\C[x_1,\dots,x_n]$, with positive integer weights $w_1\ge w_2\ge\dots\ge w_n$ of the variables. A polynomial is weighted homogeneous of weight $e$ if all its monomials $x_1^{c_1}\cdots x_n^{c_n}$ have $\sum_ic_iw_i=e$. $F\in P$ is weighted homogeneous of degree $d$ and has an isolated critical point at $0$. Throughout the paper weights and degrees are assumed to be integers, and $\gcd(w_1,\dots,w_n)=1$; since $F\in\mathfrak m^3$ (see the reduction to $d>2w_1$ below), this makes $(w_1,\dots,w_n,d)$ unique for $F$ [Sa71]. There are two exceptions: Statement~1 as quoted from [S76] below, and its discussion in Section~2.1. There, as in [S76], weights are normalized: $\alpha_i=w_i/d$ (rational; $0<\alpha_i<1/2$ because $d>2w_1$, while [S76] allows $\alpha_i\le1/2$), $F$ has degree $1$, and a weighted homogeneous polynomial of weight $e$ has normalized weight $e/d$. $F_i=\partial F/\partial x_i$, $J=J_F=(F_1,\dots,F_n)$, $\nabla F=(F_1,\dots,F_n)^T$, and $\Mil=P/J$ is the Milnor algebra, graded by weight. For any polynomial $\varphi$, $J_\varphi$ is the ideal generated by its partial derivatives. $\mathfrak m=(x_1,\dots,x_n)$.

\paragraph{Grading.} $P=\bigoplus_eP_e$, where $P_e$ is the finite-dimensional space of weighted homogeneous polynomials of weight $e$. Since $J$ is generated by weighted homogeneous polynomials, $J=\bigoplus_e(J\cap P_e)$: a polynomial lies in $J$ if and only if each of its weighted homogeneous components does. So every membership in $J$ and every relation among $F_1,\dots,F_n$ can be checked separately in each weight. In particular, in an identity $A=\sum_jc_jF_j$ with $A$ weighted homogeneous of weight $e$, each $c_j$ can be replaced by its component of weight $e-(d-w_j)$. Likewise, a vector field $D$ with $D(J)\subset J$ is a sum of weighted homogeneous vector fields, and each of them maps $J$ into $J$. So it is enough to consider derivations of a single weight $-q$.

We assume $d>2w_1$. This is no loss: by K.~Saito [Sa71], after a coordinate change $F=x_1^2+\dots+x_k^2+G$, where $G\in\mathfrak m^3$ is in the other variables, which keep their weights, and $\Mil\cong\operatorname{Mil}(G)$ as graded algebras. (This assumption excludes any quadratic terms in $F$; for example, $F=xy+z^3$ is not possible.) With $d>2w_1$, every monomial of ordinary degree $\le2$ has weight $\le2w_1<d$, so $F\in\mathfrak m^3$.

A derivation of $\Mil$ of weight $-q$, $q>0$, is induced by a vector field $D=\sum_ip_i\partial_i$ with $p_i\in P$ weighted homogeneous of weight $w_i-q$ and $D(J)\subset J$. The least weight of a nonzero element of $J$ is $d-w_1$, and $w_i-q<w_1<d-w_1$, so the $p_i$ are determined by the derivation, and the derivation is $0$ if and only if all $p_i=0$.

\paragraph{Statement 1 of [S76].} \emph{Let $F\in J$, $d_\omega F=F$, $\dim_{\C}Q(F)=\mu<\infty$, $\varphi\in M\subset\C[[x_1,\dots,x_n]]$, $i(\varphi)\subset i(F)$. Then $\varphi\in i(F)$ and $\lambda(\varphi)\ge1$.}

$J$ is the set of germs $F$ of holomorphic functions at $0\in\C^n$ with $F(0)=0$, $\operatorname{grad}F(0)=0$ (in the rest of this paper $J$ denotes the Jacobian ideal). $d_\omega$ is differentiation along the vector field $\omega=\sum_i\alpha_ix_i\partial_{x_i}$ with rational $0<\alpha_i\le1/2$; so $d_\omega F=F$ says that $F$ has degree $1$ for the normalized weights $\alpha_i=w_i/d$ (see Notation). $i(F)\subset\C[[x_1,\dots,x_n]]$ is the ideal generated by the partial derivatives of $F$, i.e.\ $J_F$, and $Q(F)=\C[[x_1,\dots,x_n]]/i(F)$; so $\mu<\infty$ says that the critical point $0$ is isolated. $M$ is the maximal ideal of $\C[[x_1,\dots,x_n]]$; so $\varphi\in M$ means $\varphi(0)=0$. $\lambda(\varphi)$ is the normalized weight of $\varphi$: if $\varphi$ is weighted homogeneous of weight $e$, then $d_\omega\varphi=(e/d)\varphi$ and $\lambda(\varphi)=e/d$; in general $\lambda(\varphi)=e/d$, where $e$ is the weight of the lowest weighted homogeneous part of $\varphi$. So $\lambda(\varphi)\ge1$ says that this lowest part has weight $\ge d$, the degree of $F$. Since $i(F)$ is generated by weighted homogeneous polynomials, membership in $i(F)$ can be checked on each weighted homogeneous part (see Grading), so for a polynomial $\varphi$ the conclusion $\varphi\in i(F)$ is $\varphi\in J_F$.

The first part is proved in [S76]. For each $F$, the second part is equivalent to: $\Mil$ has no nonzero derivation of negative weight (Section~2.1). Section~2 describes the history of the statement. Our result is the following.

\begin{theorem}\label{thm:main}
Let $F$ be as above, $n\ge3$, $q\ge1$, and
\[
M=\min(w_3-2w_n,\ w_1-3w_n).
\]
If $q>M$, or if $q>M-(w_{n-1}-w_n)$ and $w_n$ does not divide $d-q$, then $\Mil$ has no nonzero derivation of weight $-q$. In particular, if $w_3\le2w_n$ or $w_1\le3w_n$, then $\Mil$ has no nonzero derivation of negative weight, and the second part of Statement~1 holds for $F$.
\end{theorem}

Theorem~\ref{thm:main} is new, and its proof is independent of the published proofs. Among the published results whose proofs are available, it is stronger than every condition on the weights alone that applies for all $n$ in the papers listed in Section~2.4: it strictly contains the condition $w_1\le2w_n$ of [YZ16] (whose proof needs the repair described in Section~2.6) and the case of equal weights. It does not contain, and is not contained in, the condition $d>w_1^2-w_1+1$ of [CY26], which uses the degree $d$ (Examples~\ref{ex:126655} and~\ref{ex:3331}). The announced general proof [Al26] is not part of this comparison, because its proof is not available (Section~2.5).

Theorem~\ref{thm:main} follows from a more flexible criterion, Theorem~\ref{thm:rows}: for each row of the matrix $C$ of Section~3 one of three conditions on the weights rules out that this row is the first nonzero row of $C$, and different rows may be excluded by different conditions (Example~\ref{ex:77732}).

\section{History of the statement}

\subsection{Statement 1 and the paper [S76]}

Statement~1 appears in the author's paper [S76] of 1976. The first part, $\varphi\in J_F$, is Lemma~2 of [S76] (pp.~59--60), and its proof is correct: $J_F$ is generated by weighted homogeneous polynomials, so each weighted homogeneous part $\varphi_e$ of $\varphi$ (of weight $e$; here $e>0$, because $\varphi\in M$ gives $\varphi(0)=0$) also satisfies $J_{\varphi_e}\subset J_F$, and Euler's identity $\varphi_e=e^{-1}\sum_kw_kx_k\,\partial_k\varphi_e$ puts $\varphi_e$ in $J_F$.

For the second part [S76] gives a residue argument (p.~61), and there is a gap in that argument. Let $\varphi$ be weighted homogeneous of normalized weight $\lambda<1$ (see Notation); here $\lambda>0$, because $\varphi\in M$ gives $\varphi(0)=0$. The argument substitutes $x_i=t^{\alpha_i}x_i'$ in the residue integral
\[
I=\int\frac{\det(\partial_i\partial_jF)\,dx}{\prod_i(\partial_i\varphi+F_i)}
\]
and writes, in formula~(5) of [S76], $I=t^{k(1-\lambda)}\Pi_t$, where $k$ is the number of $i$ with $\partial_i\varphi\ne0$ and $\Pi_t$ is the same integral with denominators $\partial_i\varphi+t^{1-\lambda}F_i$. Comparing the two sides it concludes $k=n$, which contradicts $\partial_1\varphi=0$ (the least weight in $J_F$ is $d-w_1$). But the substitution multiplies every factor $\partial_i\varphi+F_i$, also those with $\partial_i\varphi=0$, by $t^{\lambda-\alpha_i}$, so the correct factor is $t^{n(1-\lambda)}$, not $t^{k(1-\lambda)}$. With the correct factor, formula~(5) is the residue identity for $\varphi+t^{1-\lambda}F$, which holds for every $\varphi$, and $k=n$ does not follow. So the second part is not proved in [S76].

\paragraph{The second part in terms of derivations.} For each $F$ the second part is equivalent to: $\Mil$ has no nonzero derivation of negative weight. If $J_\varphi\subset J_F$ with $\varphi$ of weight $<d$ (and of weight $>0$, because $\varphi\in M$ gives $\varphi(0)=0$), then $\varphi=D(F)$ for a vector field $D$ of negative weight with $D(J_F)\subset J_F$. Indeed, $\varphi\in J_F$ by quasihomogeneity (Euler's identity, as in the first part), so $\varphi=\sum_jp_j\partial_jF=D(F)$ with $D=\sum_jp_j\partial_j$. Conversely $D(F)$ satisfies $J_{D(F)}\subset J_F$, and $D\ne0$ gives $D(F)\ne0$ by Lemma~\ref{lem:A}.

\subsection{Halperin's conjecture}

The first publication of S.~Halperin's conjecture in rational homotopy theory was in 1977 [Ha77]: for every orientable fibration whose fibre $X$ is an elliptic space with cohomology only in even degrees, the rational Serre spectral sequence collapses at $E_2$. For such $X$, $H^*(X;\mathbb Q)$ is a graded Artinian complete intersection $\mathbb Q[x_1,\dots,x_n]/(f_1,\dots,f_n)$ ([Ha77], 1977). In 1982 W.~Meier [Me82, Theorem~A] showed that the collapse for every orientable fibration with fibre $X$ is equivalent to $H^*(X;\mathbb Q)$ having no nonzero derivation of negative degree. Since 1982 Halperin's conjecture has been stated in this algebraic form: a graded Artinian complete intersection has no nonzero derivation of negative weight ([CY26, Conjecture~1.1], 2026). One has: $\Mil=P/(\partial F/\partial x_1,\dots,\partial F/\partial x_n)$ is a graded Artinian complete intersection. So the second part of Statement~1 would follow from Halperin's conjecture for general graded Artinian complete intersections.

\subsection{The name used since 1995}

The case of Milnor algebras was proposed independently by S.\,S.-T.~Yau, and in the literature since 1995 it has been called Yau's conjecture ([CXY95], [Ch95]; [CY26], Conjecture~1.2). It is the second part of Statement~1 of [S76]. In these terms it was proved for $n=2,3$ by H.~Chen, Y.-J.~Xu and Yau [CXY95], for $n=4$ by H.~Chen [Ch95], for $w_n\ge w_1/2$ by Yau and H.~Zuo [YZ16], and for $d>w_1^2-w_1+1$ by B.~Chen and Yau [CY26]. The proof in [YZ16] has a gap (Section~2.6).

\subsection{Published work}

The tables list the published papers on the second part of Shoshitaishvili's Statement~1, in the form for Milnor algebras and in the general form for Artinian complete intersections (Halperin's conjecture), which contains it as the case $f_i=\partial F/\partial x_i$. The lists are compiled from the bibliographies of [CY26] and [CYZ19] and the papers cited there.

\medskip
\noindent\emph{Milnor algebras.}

\smallskip
\noindent{\small\begin{tabular}{>{\raggedright\arraybackslash}p{0.30\textwidth}p{0.07\textwidth}>{\raggedright\arraybackslash}p{0.55\textwidth}}
\toprule
authors & year & content\\
\midrule
A.\,N.~Shoshitaishvili [S76] & 1976 & Statement~1: $\varphi(0)=0$ and $J_\varphi\subset J_F$ imply $\varphi\in J_F$ (proved) and weighted degree $\ge d$ (residue argument with a gap)\\
H.~Chen, Y.-J.~Xu, S.\,S.-T.~Yau [CXY95] & 1995 & the statement in terms of derivations; proof for $n=2,3$\\
H.~Chen [Ch95] & 1995 & proof for $n=4$\\
Y.-J.~Xu, S.\,S.-T.~Yau [XY96] & 1996 & micro-local characterization of quasi-homogeneous singularities by the derivation Lie algebra of $\Mil$, assuming the statement\\
S.\,S.-T.~Yau, H.~Zuo [YZ16] & 2016 & $w_n\ge w_1/2$; the proof has a gap (Section~2.6)\\
B.~Chen, S.\,S.-T.~Yau [CY26] & 2026 & $d>w_1^2-w_1+1$\\
\bottomrule
\end{tabular}}

\medskip
\noindent\emph{Artinian complete intersections: Halperin's conjecture.}

\smallskip
\noindent{\small\begin{tabular}{>{\raggedright\arraybackslash}p{0.30\textwidth}p{0.07\textwidth}>{\raggedright\arraybackslash}p{0.55\textwidth}}
\toprule
authors & year & content\\
\midrule
S.~Halperin [Ha77], [FHT01] & 1977 & the conjecture (collapse of the rational Serre spectral sequence), first published in [Ha77]\\
J.-C.~Thomas [Th81] & 1981 & $n=2$\\
W.~Meier [Me82] & 1982 & equivalence with the absence of negative derivations; homogeneous fibres\\
H.~Shiga, M.~Tezuka [ST87] & 1987 & cohomology of homogeneous spaces $G/U$ of equal rank\\
G.~Lupton [Lu90] & 1990 & $n=3$\\
M.~Markl [Mar90] & 1990 & relations homogeneous in the ordinary degree\\
A.\,G.~Aleksandrov, B.~Martin [AM92] & 1992 & equal weights; $n=2$\\
S.~Papadima, L.~Paunescu [PP96] & 1996 & $P/(f_2,\dots,f_n)$ reduced and $\deg f_1\ge\deg f_i$\\
H.~Chen [Ch99] & 1999 & $n=3$\\
H.~Chen, S.\,S.-T.~Yau, H.~Zuo [CYZ19] & 2019 & graded Artinian algebras (not necessarily complete intersections) whose relations have large degrees\\
L.~Kennard, Y.~Wu [KW23] & 2023 & formal dimension $\le20$\\
A.\,G.~Aleksandrov [Al26] & 2026 & announces a proof of the conjecture in general (Section~2.5)\\
\bottomrule
\end{tabular}}

\smallskip
Conjectures of the same type have been proposed for other algebras attached to $F$: the $k$-th Hessian algebras [MYZ23] and the higher Nash blow-up local algebras [HMYZ23].

\subsection{The announced general proof}

A.\,G.~Aleksandrov, \emph{Derivations of Artinian complete intersections} [Al26] (Algebra i Analiz 38:1 (2026), in Russian; received 9 June 2025), states in its abstract: ``The main result is a proof of a long-standing conjecture which asserts that there are no derivations of negative weights on positively graded Artinian complete intersections.'' The same result is announced in A.\,G.~Aleksandrov, \emph{Vector fields on zero-dimensional complete intersections} [Al25], and was conjectured in [Al91]. Both abstracts name as the key step that the canonical primary decomposition of an ideal invariant under a derivation is itself invariant. For weighted homogeneous $F$ with an isolated critical point, $\Mil$ is a positively graded Artinian complete intersection, so the announced theorem contains the second part of Statement~1 for every such $F$.

The full texts of [Al26] and [Al25] are not openly available, and we could not read the proof. B.~Chen and Yau [CY26], published in 2026, list the general case as open. The results of the present paper do not use [Al26] or [Al25].

\subsection{The gap in [YZ16]}

Let $D$ be a vector field of weight $-q<0$ with $D(J_F)\subset J_F$. Lemma~3.1 of [YZ16] gives, for $i\ge2$,
\[
\partial_i(DF)=q_1F_1+\dots+q_{i-1}F_{i-1},
\]
with $q_j$ weighted homogeneous of weight $w_j-w_i-q$ and depending only on $x_{j+1},\dots,x_n$. The proof of the Main Theorem of [YZ16] (p.~24) then says: ``If $q(x_2,x_3,\dots,x_n)$ is not $0$, then $\deg q(x_2,x_3,\dots,x_n)\ge\operatorname{wt}(x_n)=w_n$'', and concludes that all $q_j=0$. A nonzero constant $q_j$ has weight $0$, not $\ge w_n$, and the weights allow a nonzero constant $q_j$ exactly when $w_j=w_i+q$. Excluding constants needs the isolatedness of the critical point, which this step does not use. The step fails for a non-isolated $F$ that satisfies the other hypotheses: $F=zy^3$ with $\operatorname{wt}(z)=2$, $\operatorname{wt}(y)=1$, $d=5$, and $D=y\,\partial_z$ (weight $-1$) has $D(J_F)\subset J_F$ and $\partial_y(DF)=4y^3=4\,\partial F/\partial z$, so $q_1=4\ne0$, while $D\ne0$. For isolated $F$ the constant case is settled by Lemma~\ref{lem:C}(a) below, and the Main Theorem of [YZ16] ($w_1\le2w_n$) is contained in Theorem~\ref{thm:main} ($w_1\le3w_n$).

\section{The matrix $C$ and its first nonzero row}

\begin{lemma}\label{lem:A}
Let $q>0$ be arbitrary. If $D$ induces a derivation of weight $-q<0$ and $D(F)=0$, then $D=0$.
\end{lemma}

\begin{proof}
$D(F)=\sum_ip_iF_i=0$ is a relation among $F_1,\dots,F_n$, which form a regular sequence (the critical point is isolated). The relations of a regular sequence are generated by the trivial ones $F_je_i-F_ie_j$, so every $p_i$ lies in $J$. But $p_i$ has weight $w_i-q<d-w_1$, the least weight of a nonzero element of $J$. So $p_i=0$.
\end{proof}

\paragraph{The matrix $C$.} Let $D=\sum_ip_i\partial_i$ induce a derivation of weight $-q<0$, and $g=D(F)$. Since $D(F_i)\in J$, there is a matrix $B$ of weighted homogeneous polynomials with $D(\nabla F)=B\nabla F$, $B_{ij}$ of weight $w_j-w_i-q$. Let $P_D=(\partial p_i/\partial x_j)$. From $g=\sum_kp_kF_k$,
\[
\partial_ig=\sum_k(\partial_ip_k)F_k+\sum_kp_k\,\partial_kF_i=(P_D^T\nabla F)_i+(D(\nabla F))_i,
\]
so
\[
\nabla g=C\,\nabla F,\qquad C=B+P_D^T,
\]
an identity of polynomials, with $C_{ij}$ weighted homogeneous of weight $w_j-w_i-q$. $C_{ij}\ne0$ needs $w_j\ge w_i+q>w_i$, hence $j<i$: $C$ is strictly lower triangular.

If $C=0$, then $\nabla g=0$, $g$ is a constant of weight $d-q>0$, so $g=0$ and $D=0$ by Lemma~\ref{lem:A}. So to prove that there is no nonzero derivation of weight $-q$ it is enough to show $C=0$.

\paragraph{The first nonzero row.} Suppose $C\ne0$, and let $k$ be its first nonzero row. Rows $1,\dots,k-1$ vanish, so $\partial_1g=\dots=\partial_{k-1}g=0$: $g$ depends only on $x_k,\dots,x_n$. Put
\[
V=\sum_{j<k}C_{kj}\,\partial_j,\qquad h=V(F)=\sum_{j<k}C_{kj}F_j=\partial_kg .
\]
Since $h=\partial_kg$ is the $k$-th row of $\nabla g=C\nabla F$ and $g$ depends only on $x_k,\dots,x_n$, $h$ also depends only on $x_k,\dots,x_n$. $V(x_j)=C_{kj}$ for $j<k$, and $V(x_i)=0$ for $i\ge k$. We call $x_j$ \emph{moved} if $V(x_j)\ne0$ and \emph{fixed} otherwise. In particular $x_k,\dots,x_n$ are fixed, so $V(h)=0$. Each term $C_{kj}\partial_j$ has weight $(w_j-w_k-q)-w_j=-(w_k+q)<0$, so $V$ lowers weight and is locally nilpotent, i.e.\ nilpotent in application to any particular polynomial. $C_{kj}$ has weight $w_j-w_k-q<w_j$, so it does not involve $x_j$ or any variable of weight $\ge w_j$.

\section{Lemmas on the first nonzero row}

\begin{lemma}\label{lem:C}
In the setting above:
\begin{itemize}
\item[(a)] every entry of row $k$ of $C$ that is a constant is $0$;
\item[(b)] row $k$ has at least two nonzero entries.
\end{itemize}
\end{lemma}

\begin{proof}
(a) Let $\Phi_\tau=\exp(\tau V)$, a polynomial automorphism of $\C^n$ for each $\tau\in\C$. Since for $h=V(F)$ one has $V(h)=0$,
\[
F\circ\Phi_\tau=F+\tau h .
\]
$F\in\mathfrak m^3$ gives $F_j\in\mathfrak m^2$, so $h\in\mathfrak m^2$, and $0$ is a critical point of $F+\tau h$. By the chain rule, $0=D\Phi_\tau(0)^T\,\nabla F(\Phi_\tau(0))$ with $D\Phi_\tau(0)$ invertible (indeed, $\Phi_{-\tau}\circ\Phi_\tau$ is the identity, so by the chain rule $D\Phi_{-\tau}(\Phi_\tau(0))\,D\Phi_\tau(0)=I$), so $\Phi_\tau(0)$ is a critical point of $F$. The critical set of $F$ is invariant under $x_i\mapsto t^{w_i}x_i$, so the isolated critical point $0$ is the only one. Hence $\Phi_\tau(0)=0$ for all $\tau$, and $V(0)=\frac{d}{d\tau}\Phi_\tau(0)|_{\tau=0}=0$, i.e.\ $C_{kj}(0)=0$ for all $j$. A weighted homogeneous constant $C_{kj}$ equals $C_{kj}(0)$, so it is $0$.

(b) Suppose row $k$ has only one nonzero entry $C_{kj}$, $j<k$. Then $h=C_{kj}F_j$ is a single term. $h$ does not contain $x_j$, since it depends only on $x_k,\dots,x_n$, and $C_{kj}$ does not contain $x_j$, since its weight is $<w_j$. With one term there is nothing to cancel the $x_j$ in $F_j$, so $F_j$ does not contain $x_j$: $0=\partial_jh=C_{kj}\,\partial_j^2F$, hence $\partial_j^2F=0$. Then $F=A+x_jB$ with $A,B$ free of $x_j$, $A\in\mathfrak m^3$, $B\in\mathfrak m^2$, and the whole $x_j$-axis is critical (there $\partial_jF=B(0)=0$ and $\partial_iF=\partial_iA(0)+x_j\,\partial_iB(0)=0$), which contradicts the isolated critical point.
\end{proof}

\begin{lemma}\label{lem:D}
Let $F\in\C[u,v,z_1,\dots,z_s]$ lie in $\mathfrak m^3$ and have an isolated critical point at $0$. Let $a\in\C[v,z]$, $b\in\C[z]$, $V=a\,\partial_u+b\,\partial_v$, and suppose $V(F)\in\C[z]$. Then $a=b=0$.
\end{lemma}

No weights are used here.

\begin{proof}
Suppose $V\ne0$ and put $h=V(F)\in\C[z]$.

\emph{Step 1: common factors.} Let $e=\gcd(a,b)$ in $\C[v,z]$ ($e=a$ if $b=0$; if $b\ne0$, $e$ divides $b$, so $e\in\C[z]$). Write $a=ea'$, $b=eb'$, $V'=a'\partial_u+b'\partial_v$. Then $e\,V'(F)=h$. If $h\ne0$: when two nonzero polynomials are multiplied, the terms with the highest powers of $u$ and $v$ multiply and cannot cancel; $h$ contains no $u$ or $v$, so neither does $V'(F)$, i.e.\ $V'(F)\in\C[z]$. If $h=0$, then $V'(F)=0$. So $V'$ satisfies the hypotheses, with $\gcd(a',b')=1$ (and $a'=1$ if $b'=0$).

\emph{Step 2: $V'(0)=0$.} $V'$ is locally nilpotent, so $\Phi_\tau=\exp(\tau V')$ is a polynomial automorphism and $F\circ\Phi_\tau=F+\tau h'$ with $h'=V'(F)\in\C[z]\cap\mathfrak m^2$. As in Lemma~\ref{lem:C}(a), $\Phi_\tau(0)$ is a critical point of $F$ for every $\tau$, depends polynomially on $\tau$ and equals $0$ at $\tau=0$; so $\Phi_\tau(0)=0$ and $(a'(0),b'(0))=0$. In particular $a'\ne1$, so $b'\ne0$ and $b'(0)=0$.

From now on write $a,b$ for $a',b'$: $\gcd(a,b)=1$, $b\in\C[z]$, $b\ne0$, $b(0)=0$, $V(F)=h\in\C[z]$.

\emph{Step 3: $F_{uu}=bK$, $F_{uv}=-aK$.} $\partial_u$ commutes with $V$ and $\partial_uh=0$, so $aF_{uu}+bF_{uv}=0$. As $a,b$ are coprime in $\C[u,v,z]$, $b$ divides $F_{uu}$, so $F_{uu}=bK$ for a polynomial $K$, and then $F_{uv}=-aK$.

\emph{Step 4: every prime factor of $b$ divides $K$.} Let $\pi\in\C[z]$ be an irreducible factor of $b$. Then $R=\C[u,v,z]/(\pi)=(\C[z]/(\pi))[u,v]$ is a domain and $\pi\nmid a$. Modulo $\pi$: $F_{uu}\equiv0$, so the class $\bar F_u$ does not involve $u$; and $V(F)\equiv\bar a\bar F_u$, so $\bar a\,\bar F_u=\bar h$ with $\bar a\ne0$ and $\bar h\in\C[z]/(\pi)$. In $(\C[z]/(\pi))[v]$ the product of two nonzero elements is nonzero, so if $\bar F_u$ contained $v$, the highest power of $v$ in $\bar a$ times the highest power in $\bar F_u$ would appear in the product, and $\bar h$ would contain $v$. Hence $\bar F_u$ contains no $v$ (if $\bar h=0$, then $\bar F_u=0$ for the same reason). So $\partial_v\bar F_u=0$; since $\pi$ does not involve $v$, this means $\partial_vF_u=F_{uv}=-aK\equiv0\pmod\pi$, and $\pi\mid K$ because $\pi$ does not divide $a$.

\emph{Step 5.} By Step~4, $\operatorname{rad}(b)\mid K$, so $F_{uu}=b\operatorname{rad}(b)L$ and $F_{uv}=-a\operatorname{rad}(b)L$ for a polynomial $L$. Since $b(0)=0$, $b$ and $\operatorname{rad}(b)$ lie in $(z_1,\dots,z_s)$. Hence $F_{uu}\in(z)^2$ and $F_{uv}\in(z)$.

\emph{Step 6: the $u$-axis is critical.} On the $u$-axis ($v=0$, $z=0$): $F_{uu}=0$, so $F_u=F_u(0)=0$; $F_{uv}=0$, so $F_v=F_v(0)=0$; and $\partial_{z_i}F_{uu}\in(z)$, so $F_{uuz_i}=0$ there and $F_{z_i}(u,0,0)=F_{z_i}(0)+u\,F_{uz_i}(0)=0$, since $F\in\mathfrak m^3$. So $\nabla F$ vanishes on the $u$-axis, contradicting isolatedness.
\end{proof}

\begin{corollary}\label{cor:three}
The first nonzero row of $C$ has at least three nonzero entries, and all of them are nonconstant.
\end{corollary}

\begin{proof}
By Lemma~\ref{lem:C} the nonzero entries are nonconstant and there are at least two. Suppose there are exactly two, $C_{kj_1},C_{kj_2}$ with $j_1<j_2<k$. Put $u=x_{j_1}$, $v=x_{j_2}$, and let $z$ be the other variables. $C_{kj_2}$ has weight $<w_{j_2}\le w_{j_1}$, so it involves neither $u$ nor $v$; $C_{kj_1}$ does not involve $u$. $V=C_{kj_1}\partial_u+C_{kj_2}\partial_v$ has $V(F)=h\in\C[z]$. Lemma~\ref{lem:D} gives $V=0$, a contradiction.
\end{proof}

\section{Linear vector fields}

\begin{lemma}\label{lem:G}
Let $R$ be a commutative ring containing $\tfrac12$, $a=(a_1,\dots,a_m)^T$ with $a_1,\dots,a_m$ a regular sequence in $R$, $I=(a_1,\dots,a_m)$, and $H$ a symmetric $m\times m$ matrix over $R$ with $Ha=0$. Then every entry of $H$ lies in $I^2$.
\end{lemma}

\begin{proof}
For a regular sequence, $I/I^2$ is a free $R/I$-module with basis $\bar a_1,\dots,\bar a_m$, and $I^2/I^3$ is a free $R/I$-module with basis $\bar a_j\bar a_l$ ($j\le l$) ([Mat86], \S16: a regular sequence is quasi-regular). From $\sum_jH_{ij}a_j=0$ and freeness of $I/I^2$, $H_{ij}\in I$; write $H_{ij}=\sum_lT_{ijl}a_l$. Symmetry of $H$ gives $\sum_l(T_{ijl}-T_{jil})a_l=0$, so $T_{ijl}\equiv T_{jil}\pmod I$. Next, $\sum_{j,l}T_{ijl}a_ja_l=0$, so in $I^2/I^3$ the coefficient of $\bar a_j\bar a_l$ vanishes: $\bar T_{ijl}=-\bar T_{ilj}$ (using $\tfrac12$ for $j=l$). So $\bar T$ is symmetric in the first two indices and antisymmetric in the last two, and
\[
\bar T_{ijl}=\bar T_{jil}=-\bar T_{jli}=-\bar T_{lji}=\bar T_{lij}=\bar T_{ilj}=-\bar T_{ijl},
\]
so $\bar T=0$, $T_{ijl}\in I$ and $H_{ij}\in I^2$.
\end{proof}

\paragraph{The $\mathfrak{sl}_2$ grading.} Let $N\ne0$ be a nilpotent linear vector field on $\C^n$. $N$ acts on linear forms as a nilpotent operator; choose a Jordan basis $y_{b,1},\dots,y_{b,s_b}$ (blocks $b$) with $N(y_{b,i})=y_{b,i+1}$, $N(y_{b,s_b})=0$, and use it as coordinates. There are linear vector fields $E$, $H$ with
\[
[H,E]=2E,\qquad [H,N]=-2N,\qquad [E,N]=H,\qquad H(y_{b,i})=(s_b+1-2i)\,y_{b,i},
\]
that is, $(E,H,N)$ is an $\mathfrak{sl}_2$-triple (this is the Jacobson--Morozov theorem [Mo42], [Ja51]). Here $\mathfrak{sl}_2$ is the Lie algebra of complex $2\times2$ matrices with trace $0$; its basis $e=\begin{pmatrix}0&1\\0&0\end{pmatrix}$, $h=\begin{pmatrix}1&0\\0&-1\end{pmatrix}$, $f=\begin{pmatrix}0&0\\1&0\end{pmatrix}$ satisfies the same relations, so $e\mapsto E$, $h\mapsto H$, $f\mapsto N$ is an action of $\mathfrak{sl}_2$ on polynomials. We need it for (S1) and (S2) below, which are facts about finite-dimensional $\mathfrak{sl}_2$-modules. We call $s_b+1-2i$ the $H$-weight of $y_{b,i}$ and extend it additively to monomials. Every space of homogeneous polynomials of a fixed ordinary degree is a finite-dimensional $\mathfrak{sl}_2$-module, and $N$ lowers the $H$-weight by $2$. Two standard facts:
\begin{itemize}
\item[(S1)] A nonzero $H$-homogeneous polynomial $f$ with $N(f)=0$ has $H$-weight $\le0$.
\item[(S2)] For any polynomial $G$, the $H$-weight-$0$ part of $N(G)$, if it is killed by $N$, is $0$.
\end{itemize}

\begin{lemma}\label{lem:J}
Let $F\in\mathfrak m^3$ have an isolated critical point at $0$, and let $N\ne0$ be a nilpotent linear vector field. It is not possible that $N(F)=h$ where, in some linear coordinates, $h$ contains only coordinates not moved by $N$ (coordinates $x_i$ with $N(x_i)=0$).
\end{lemma}

\begin{proof}
Use the Jordan coordinates. The linear forms $\ell$ with $N(\ell)=0$ are the combinations of the $y_{b,s_b}$, so $h$ is a polynomial in the $y_{b,s_b}$, and $N(h)=0$.

\emph{Step 1: $F$ has $H$-weights $\le1$.} Split $h=\sum_\mu h_\mu$ and $F=\sum_\mu F_\mu$ into $H$-homogeneous parts. Each $h_\mu$ is killed by $N$, so $\mu\le0$ by (S1), and $h_0=0$ by (S2). For $\mu\ge2$, $N(F_\mu)=h_{\mu-2}=0$, and $F_\mu$ has positive weight, so $F_\mu=0$ by (S1).

\emph{Step 2: a critical line.} Let $s\ge2$ be the largest block size, $b_0$ a block of size $s$, $U=y_{b_0,1}$ (of $H$-weight $s-1$), and $L$ the $U$-axis. A monomial of $F$ contributes to a first derivative of $F$ at points of $L$ only if it is $U^a$ or $U^ax$ with $x\ne U$ a coordinate; as $F\in\mathfrak m^3$, only $U^a$ with $a\ge3$ and $U^ax$ with $a\ge2$ occur. Every coordinate has $H$-weight $\ge-(s-1)$, so these monomials have weight $\ge3(s-1)$, resp.\ $\ge s-1$.

If $s\ge3$, all these weights are $\ge2$, so by Step~1 none of these monomials occurs, and $\nabla F=0$ on $L$.

If $s=2$, the only remaining candidates are $U^2y$ with $y$ of weight $-1$, the second coordinate of a block of size $2$. Write $t_1,\dots,t_p$ for the first and $b_1,\dots,b_p$ for the second coordinates of the blocks of size $2$ ($t_1=U$), and $z$ for the coordinates of the blocks of size $1$; then $N=\sum_ib_i\partial_{t_i}$ lowers the degree in $t$ by $1$, and $h\in\C[b,z]$. Let $F^{(2)}$ be the part of $F$ of degree $2$ in $t$; $N(F^{(2)})$ is the part of $h$ of degree $1$ in $t$, which is $0$. The matrix $\mathcal H=(\partial_{t_i}\partial_{t_l}F^{(2)})$ has entries in $\C[b,z]$ and $\sum_ib_i\mathcal H_{il}=\partial_{t_l}N(F^{(2)})=0$. As $b_1,\dots,b_p$ is a regular sequence in $\C[b,z]$, Lemma~\ref{lem:G} gives $\mathcal H_{11}\in(b)^2$. The coefficient of $U^2$ in $F$ is $\tfrac12\mathcal H_{11}$, so $F$ has no monomial $U^2b_j$. Again $\nabla F=0$ on $L$, contradicting isolatedness.
\end{proof}

\section{Proof of Theorem~\ref{thm:main}}

\paragraph{Notation.} $S=\{m_1w_1+\dots+m_nw_n:\ m_i\in\mathbb Z_{\ge0}\}$ is the set of weights of monomials, and $S_k=\{m_kw_k+\dots+m_nw_n:\ m_i\in\mathbb Z_{\ge0}\}$ the set of weights of monomials in $x_k,\dots,x_n$. For $q>0$ and $2\le k\le n$,
\[
I_{k,q}=\{\,j<k:\ w_j-w_k-q>0,\ w_j-w_k-q\in S\,\},
\]
the columns in which row $k$ of $C$ can have a nonconstant entry.

\begin{theorem}\label{thm:rows}
Let $F$ be as in Section~1 and $q>0$. Suppose that for every $k$ with $2\le k\le n$ at least one of the following holds:
\begin{itemize}
\item[(a)] $|I_{k,q}|\le2$;
\item[(b)] $d-w_k-q\notin S_k$;
\item[(c)] for no $j\in I_{k,q}$ is $w_j-w_k-q$ the weight of a monomial of ordinary degree $\ge2$.
\end{itemize}
Then $\Mil$ has no nonzero derivation of weight $-q$.
\end{theorem}

\begin{proof}
Suppose $C\ne0$, let $k$ be its first nonzero row, and let $V$, $h$ be as in Section~3. By Corollary~\ref{cor:three}, row $k$ has at least three nonzero entries, all nonconstant, so they lie in columns of $I_{k,q}$.

If (a) holds for this $k$, this is a contradiction.

If (b) holds, then $h=\sum_{j<k}C_{kj}F_j$, a polynomial in $x_k,\dots,x_n$ of weight $d-w_k-q\notin S_k$, is $0$. This is a relation among the regular sequence $F_1,\dots,F_n$, so every $C_{kj}$ lies in $J$ (as in Lemma~\ref{lem:A}); $C_{kj}$ has weight $<w_1<d-w_1$, so $C_{kj}=0$, a contradiction.

If (c) holds, every nonzero $C_{kj}$ is a nonconstant weighted homogeneous polynomial whose weight is not that of a monomial of degree $\ge2$, so it is a linear form. Then $V$ is a nonzero linear vector field of weight $-(w_k+q)<0$, hence nilpotent, and $V(F)=h$ contains only $x_k,\dots,x_n$, which $V$ does not move. Lemma~\ref{lem:J} gives a contradiction.

So $C=0$, and there is no nonzero derivation of weight $-q$ (Section~3).
\end{proof}

Different rows may use different conditions (Example~\ref{ex:77732}).

\begin{proof}[Proof of Theorem~\ref{thm:main}]
We check the hypothesis of Theorem~\ref{thm:rows} in four cases. The least positive element of $S$ is $w_n$, and a monomial of ordinary degree $\ge2$ has weight $\ge2w_n$.

\emph{Case $q>w_3-2w_n$.} For $j\ge3$ and every $k>j$, $0<w_j-w_k-q\le w_3-w_n-q<w_n$ is impossible for an element of $S$, so $I_{k,q}\subset\{1,2\}$ and (a) holds for every $k$.

\emph{Case $q>w_1-3w_n$.} For every $j<k$, $w_j-w_k-q\le w_1-w_n-q<2w_n$, so (c) holds for every $k$.

In the other two cases $w_n$ does not divide $d-q$. Since $S_n=w_n\mathbb Z_{\ge0}$, $d-w_n-q\in S_n$ would mean that $w_n$ divides $d-q$; so (b) holds for $k=n$. For $k\le n-1$ we have $w_k\ge w_{n-1}$.

\emph{Case $q>w_3-w_{n-1}-w_n$.} For $3\le j<k\le n-1$, $w_j-w_k-q\le w_3-w_{n-1}-q<w_n$, so $j\notin I_{k,q}$; hence $I_{k,q}\subset\{1,2\}$ and (a) holds.

\emph{Case $q>w_1-w_{n-1}-2w_n$.} For $j<k\le n-1$, $w_j-w_k-q\le w_1-w_{n-1}-q<2w_n$, so (c) holds.

The first two cases are $q>M$; the last two are $q>M-(w_{n-1}-w_n)$ with $w_n\nmid d-q$, since $w_3-w_{n-1}-w_n=(w_3-2w_n)-(w_{n-1}-w_n)$ and $w_1-w_{n-1}-2w_n=(w_1-3w_n)-(w_{n-1}-w_n)$. If $w_3\le2w_n$ or $w_1\le3w_n$, then $M\le0<q$ for every $q\ge1$.
\end{proof}

\paragraph{Comparison.} If $w_1\le2w_n$, then $w_1\le3w_n$; so Theorem~\ref{thm:main} contains the Main Theorem of [YZ16] and the case of equal weights. For $n\le4$ the absence of negative derivations is known for all weights ([CXY95], [Ch95]). For $n\ge5$ the cases $w_3\le2w_n$ and $w_1\le3w_n$ are not covered by the results listed in Section~2, except for homogeneous $F$ and for $d>w_1^2-w_1+1$ [CY26]. The condition of Theorem~\ref{thm:main} and the condition $d>w_1^2-w_1+1$ do not contain each other (Examples~\ref{ex:126655} and~\ref{ex:3331}).

\section{Examples}

\begin{example}\label{ex:126655}
Weights $(12,6,6,5,5)$, $d=60$, $F=x_1^5+x_2^{10}+x_3^{10}+x_4^{12}+x_5^{12}$. Here $w_1=12>2w_5$, $F$ is not homogeneous, and $d=60\le w_1^2-w_1+1=133$. Since $w_1\le3w_5$ (and $w_3\le2w_5$), $M\le0$, and by Theorem~\ref{thm:main} $\Mil$ has no nonzero derivation of negative weight for every isolated $F$ with these weights.
\end{example}

\begin{example}\label{ex:3331}
Weights $(4,4,4,1,1)$, $d=16$, $F=x_1^4+x_2^4+x_3^4+x_4^{16}+x_5^{16}$. Here $d=16>w_1^2-w_1+1=13$, so [CY26] applies. On the other hand $M=\min(4-2,\,4-3)=1$ and $M-(w_4-w_5)=1$, so Theorem~\ref{thm:main} does not cover $q=1$.
\end{example}

\begin{example}\label{ex:77732}
Weights $(7,7,7,3,2)$, $d=42$, $F=x_1^6+x_2^6+x_3^6+x_4^{14}+x_5^{21}$; here $d=42\le w_1^2-w_1+1=43$. $M=\min(7-4,\,7-6)=1$, so $q\ge2$ is covered by $q>M$. For $q=1$: $M-(w_4-w_5)=0<1$ and $w_5=2$ does not divide $d-q=41$, so Theorem~\ref{thm:main} applies. In Theorem~\ref{thm:rows}: row~$5$ satisfies (b) only ($39\notin2\mathbb Z$), and row~$4$ satisfies (c) only ($3$ is not the weight of a monomial of degree $\ge2$). So $\Mil$ has no nonzero derivation of negative weight for every isolated $F$ with weights $(7,7,7,3,2)$ and even $d$.
\end{example}

\section*{Acknowledgments}
The author used the AI assistants Claude (Anthropic) and ChatGPT (OpenAI) for help in this work.

{\small
\begin{thebibliography}{HMYZ23}
\bibitem[Al91]{Al91} A.\,G.~Aleksandrov, Vector fields on a complete intersection, Funktsional.\ Anal.\ i Prilozhen.\ 25:4 (1991), 64--66; English transl.: Funct.\ Anal.\ Appl.\ 25:4 (1991), 283--284, doi:10.1007/bf01080081.
\bibitem[Al25]{Al25} A.\,G.~Aleksandrov, Vector fields on zero-dimensional complete intersections, Izv.\ Ross.\ Akad.\ Nauk Ser.\ Mat.\ 89:2 (2025), doi:10.4213/im9294.
\bibitem[Al26]{Al26} A.\,G.~Aleksandrov, Derivations of Artinian complete intersections, Algebra i Analiz 38:1 (2026), 1--39 (in Russian).
\bibitem[AM92]{AM92} A.\,G.~Aleksandrov, B.~Martin, Derivations and deformations of Artinian algebras, Beitr\"age Algebra Geom.\ (1992), 115--130.
\bibitem[Ch95]{Ch95} H.~Chen, On negative weight derivations of the moduli algebras of weighted homogeneous hypersurface singularities, Math.\ Ann.\ 303 (1995), 95--107, doi:10.1007/bf01460981.
\bibitem[Ch99]{Ch99} H.~Chen, Nonexistence of negative weight derivations on graded Artin algebras: a conjecture of Halperin, J.\ Algebra 216 (1999), 1--12, doi:10.1006/jabr.1998.7750.
\bibitem[CXY95]{CXY95} H.~Chen, Y.-J.~Xu, S.\,S.-T.~Yau, Nonexistence of negative weight derivation of moduli algebras of weighted homogeneous singularities, J.\ Algebra 172 (1995), 243--254, doi:10.1016/s0021-8693(05)80001-8.
\bibitem[CY26]{CY26} B.~Chen, S.\,S.-T.~Yau, On nonexistence of negative weight derivations on moduli algebras: Yau's conjecture, Sbornik: Mathematics 217:1 (2026), 105--128, doi:10.4213/sm10201e.
\bibitem[CYZ19]{CYZ19} H.~Chen, S.\,S.-T.~Yau, H.~Zuo, Non-existence of negative weight derivations on positively graded Artinian algebras, Trans.\ Amer.\ Math.\ Soc.\ 372 (2019), 2493--2535, doi:10.1090/tran/7628.
\bibitem[FHT01]{FHT01} Y.~F\'elix, S.~Halperin, J.-C.~Thomas, Rational homotopy theory, Grad.\ Texts in Math.\ 205, Springer, New York, 2001.
\bibitem[Ha77]{Ha77} S.~Halperin, Finiteness in the minimal models of Sullivan, Trans.\ Amer.\ Math.\ Soc.\ 230 (1977), 173--199, doi:10.1090/s0002-9947-1977-0461508-8.
\bibitem[HMYZ23]{HMYZ23} N.~Hussain, G.~Ma, S.\,S.-T.~Yau, H.~Zuo, Higher Nash blow-up local algebras of singularities and its derivation Lie algebras, J.\ Algebra 618 (2023), 165--194, doi:10.1016/j.jalgebra.2022.11.016.
\bibitem[Ja51]{Ja51} N.~Jacobson, Completely reducible Lie algebras of linear transformations, Proc.\ Amer.\ Math.\ Soc.\ 2 (1951), 105--113, doi:10.1090/S0002-9939-1951-0049882-5.
\bibitem[KW23]{KW23} L.~Kennard, Y.~Wu, Halperin's conjecture in formal dimensions up to 20, Comm.\ Algebra 51 (2023), 3606--3622, doi:10.1080/00927872.2023.2186705.
\bibitem[Lu90]{Lu90} G.~Lupton, Note on a conjecture of Stephen Halperin's, in: Topology and Combinatorial Group Theory, Lecture Notes in Math.\ 1440, Springer, Berlin, 1990, 148--163, doi:10.1007/bfb0084459.
\bibitem[Mar90]{Mar90} M.~Markl, Towards one conjecture on collapsing of the Serre spectral sequence, Rend.\ Circ.\ Mat.\ Palermo (2) Suppl.\ 22 (1990), 151--159.
\bibitem[Mat86]{Mat86} H.~Matsumura, Commutative Ring Theory, Cambridge Studies in Advanced Mathematics~8, Cambridge University Press, 1986, doi:10.1017/CBO9781139171762.
\bibitem[Me82]{Me82} W.~Meier, Rational universal fibrations and flag manifolds, Math.\ Ann.\ 258 (1982), 329--340, doi:10.1007/bf01450686.
\bibitem[Mo42]{Mo42} V.\,V.~Morozov, On a nilpotent element in a semi-simple Lie algebra, C.\,R.\ (Doklady) Acad.\ Sci.\ URSS (N.S.) 36 (1942), 83--86.
\bibitem[MYZ23]{MYZ23} G.~Ma, S.\,S.-T.~Yau, H.~Zuo, Nonexistence of negative weight derivations of the local $k$-th Hessian algebras associated to isolated singularities, Pacific J.\ Math.\ 323 (2023), 129--172, doi:10.2140/pjm.2023.323.129.
\bibitem[PP96]{PP96} S.~Papadima, L.~Paunescu, Reduced weighted complete intersection and derivations, J.\ Algebra 183 (1996), 595--604, doi:10.1006/jabr.1996.0234.
\bibitem[S76]{S76} A.\,N.~Shoshitaishvili, Functions with isomorphic Jacobian ideals, Funktsional.\ Anal.\ i Prilozhen.\ 10:2 (1976), 57--62; English transl.: Funct.\ Anal.\ Appl.\ 10:2 (1976), 128--133, doi:10.1007/bf01077939.
\bibitem[Sa71]{Sa71} K.~Saito, Quasihomogene isolierte Singularit\"aten von Hyperfl\"achen, Invent.\ Math.\ 14 (1971), 123--142, doi:10.1007/bf01405360.
\bibitem[ST87]{ST87} H.~Shiga, M.~Tezuka, Rational fibrations, homogeneous spaces with positive Euler characteristics and Jacobians, Ann.\ Inst.\ Fourier 37:1 (1987), 81--106, doi:10.5802/aif.1078.
\bibitem[Th81]{Th81} J.-C.~Thomas, Rational homotopy of Serre fibrations, Ann.\ Inst.\ Fourier 31:3 (1981), 71--90, doi:10.5802/aif.838.
\bibitem[XY96]{XY96} Y.-J.~Xu, S.\,S.-T.~Yau, Micro-local characterization of quasi-homogeneous singularities, Amer.\ J.\ Math.\ 118 (1996), 389--399, doi:10.1353/ajm.1996.0020.
\bibitem[YZ16]{YZ16} S.\,S.-T.~Yau, H.~Zuo, Derivations of the moduli algebras of weighted homogeneous hypersurface singularities, J.\ Algebra 457 (2016), 18--25, doi:10.1016/j.jalgebra.2016.03.003.
\end{thebibliography}
}
\end{document}
